\subsection{The $s_1$ Tracker}
\label{sec:s1_controller}

The $s_1$ tracker steers along the same mountain pass, read in $s_1$,
using two channels built from $\nabla\hat{s}_1$ and the strain eigenframe.
Unlike $\hat{D}$, the fitted $\hat{s}_1$ is concave for every field and
formation (Section~\ref{sec:surrogates}), so it never shows the upward
curvature a trench has across it and the tracker cannot steer on a fitted
Hessian.

On an OECS the trench tangent is a strain eigenvector, $\mathbf{e}_2$ on
an attracting segment and $\mathbf{e}_1$ on a repelling one, and the
tracker must choose between them at every step. On the trench the
transverse derivative of $s_1$ vanishes, so $\nabla\hat{s}_1$ lies nearly
along the trench away from a stationary point. The tracker selects the
eigenvector most nearly parallel to $\nabla\hat{s}_1$ and signs it against
a reference,
\begin{equation}
    \mathbf{t}_{\mathrm{raw}} =
    \arg\max_{\mathbf{e} \in \{\mathbf{e}_1, \mathbf{e}_2\}}
    \bigl\lvert \nabla\hat{s}_1^\top \mathbf{e} \bigr\rvert,
    \qquad
    \mathbf{t} = \mathrm{sgn}\bigl(\mathbf{t}_{\mathrm{ref}}^\top
    \mathbf{t}_{\mathrm{raw}}\bigr)\,\mathbf{t}_{\mathrm{raw}},
    \label{eq:tangent_select}
\end{equation}
where $\mathbf{t}_{\mathrm{ref}}$ is the previous step's tangent
(Fig.~\ref{fig:s1_channels}(a)). On a trench that is not an OECS, the rule
takes the nearer eigenvector. The first ride step has no previous tangent,
so the measured flow $\hat{\mathbf{v}}_0$ serves as
$\mathbf{t}_{\mathrm{ref}}$. Because the measured flow is not objective, it
seeds the ride direction only once.

Behavior at the crest depends on whether the strain vanishes there. In
incompressible flow $s_1 = -r \leq 0$ reaches zero only at an isotropic
point, where $\mathbf{S}$ vanishes. Such points are in general isolated
peaks, so the crest is a smooth saddle with $r > 0$, where
$\nabla\hat{s}_1$ vanishes and tangent selection is least reliable, as it
is near a minimum. When the shear strain vanishes identically, however,
the isotropic points form curves, and the crest can sit on one at a kink
of $s_1$. The tangent changes eigenvector only at an isotropic point, and
the rule follows the swap across the kink with no dedicated test because
$\nabla\hat{s}_1$ stays large on either side. On the benchmark,
$\mathbf{S} = \mathrm{diag}(u_x, -u_x)$ has its eigenvectors on the
coordinate axes wherever $u_x \neq 0$. On the upper segment of the
separatrix $u_x < 0$, so the trench follows $\mathbf{e}_2$ and is an
attracting OECS; on the lower segment $u_x > 0$, so it follows
$\mathbf{e}_1$ and is a repelling OECS, meeting at the isotropic point at
the origin.

From the first step, the command pairs a constant-speed ride along
$\mathbf{t}$ with projected gradient descent across the trench,
\begin{equation}
    \dot{\mathbf{p}}_c = k\bigl[
    c_{\max}\tanh(1)\,\mathbf{t}
    - \mathrm{sat}\!\bigl(g_\perp\,\mathbf{P}\,\nabla\hat{s}_1\bigr)\bigr],
    \quad
    \mathbf{P} = \mathbf{I} - \mathbf{t}\mathbf{t}^\top,
    \label{eq:s1_law}
\end{equation}
with transverse gain $g_\perp$, where $\mathrm{sat}$ now saturates a
vector's magnitude at $c_{\max}$ and preserves its direction. The projector
$\mathbf{P}$ removes the along-trench component of the gradient
(Fig.~\ref{fig:s1_channels}(b)), so the cluster descends only across the
trench while riding along it at speed $k\,c_{\max}\tanh 1$.

The tracker stops at a minimum by a gradient test. The concave fit cannot
certify a minimum, but the fitted gradient tracks the true gradient even
where the fitted curvature has the wrong sign. Each end of the trench is
a true two-dimensional minimum of $s_1$, so a gradient that is small
against the depth suffices. Capture requires $\hat{s}_1 < 0$ and
\begin{equation}
    q = \frac{\rho\,\|\nabla\hat{s}_1\|}{|\hat{s}_1|} < q_{\mathrm{cap}},
    \label{eq:s1_capture}
\end{equation}
where $q$ is the fraction by which $\hat{s}_1$ changes across the
footprint. The ratio is dimensionless, so $q_{\mathrm{cap}}$ needs no
knowledge of the field's strain rate. It is large along the ride and
wherever $\hat{s}_1$ nears zero, as at an isotropic point or a gyre
center, and small only near a minimum. The tracker then drops the ride
term and descends $\hat{s}_1$, $\dot{\mathbf{p}}_c =
-k\,\mathrm{sat}(g_\perp\nabla\hat{s}_1)$, for the rest of the run. If
noise trips the test early on the trench, the descent still carries the
cluster down the trench to the minimum.

\begin{figure}[!tb]
    \centering
    \includegraphics[width=\columnwidth]{figures/s1_channels.png}
    \caption{The two channels of the $s_1$ tracker, on the benchmark
    trench at $x = 0$. (a) Tangent selection. At each sample the strain eigenframe is
    drawn, the selected eigenvector heavy and the rejected one light,
    with $\nabla\hat{s}_1$ in orange. The argmax takes
    $\mathbf{e}_2$ on the attracting segment and $\mathbf{e}_1$ on the
    repelling segment, swapping through the isotropic point at the origin.
    (b) The command. Travel along $\mathbf{t}$ is open loop at cruise
    speed, while $\mathbf{P}$ removes the along-trench component of the
    gradient so descent acts only across the trench.}
    \label{fig:s1_channels}
\end{figure}

The across-trench channel returns the cluster to the trench by the same
argument as for the $D$ tracker. It reduces to
$\dot{n} = -k\,\mathrm{sat}(a_\perp n)$, with $\kappa_\perp$ now the
transverse curvature of $s_1$ and $a_\perp = g_\perp\kappa_\perp$. No
fitted curvature enters, so $a_\perp > 0$ wherever the trench curves
upward across it, including at a smooth crest. The restoring term is lost
only where $\kappa_\perp$ vanishes, as where a curve of isotropic points
crosses the trench. There the ride carries the cluster across on the
tangent alone.